\documentclass[10pt,a4paper]{article}
\usepackage{amsmath,amssymb,amsthm,graphicx}
\title{Influential variables in completely bounded block-multilinear forms}
\author{Nikhil Bansal, Makrand Sinha, Ronald de Wolf}
\date{doi:10.4230/LIPIcs.CCC.2022.28; CC BY4.0}
\begin{document}
\maketitle
\noindent\textit{Editorial note.} Original authors: Nikhil Bansal, Makrand Sinha and Ronald de Wolf. Official CCC2022 article28 (2022-07-11), DOI10.4230/LIPIcs.CCC.2022.28. Creative Commons Attribution4.0.\par\medskip
\noindent\textit{Editorial note.} License: https://creativecommons.org/licenses/by/4.0/\par\medskip
\noindent\textit{Editorial note.} Faithful editable TeX transcription of author Theorem3 and complete root-influence/nonhomogeneous proof, necessary block-form/Fourier/completely-bounded/free-probability definitions, complete local polar-decomposition support and AppendixA.2 Lemma9 proof, exact cited standard norm/convergence results, original Corollary5 and its author derivation, and complete bibliography. Original mathematical expressions including literal Equation7 retained; no mathematical correction or new proof. Standard prior Theorems8/10 not recursively reproved. Non-exclusive arXiv version is an identity alias only, no extra preprint content published.\par\medskip
\noindent\textit{Allowed standard prerequisites (editorial).}\begin{itemize}
\item Strong asymptotic freeness of Haar unitaries. Exact author Theorem8 citing Voiculescu[35] and Collins-Male[16] retained.
\item Kemp-Speicher strong Haagerup inequality. Exact prior Theorem10[23] retained as a standard premise; all new local extensions and polar-decomposition proofs included.
\item Aaronson-Ambainis variance-reduction decision-tree theorem. Author cites [1, Theorem3.3] for the grouped simulation consequence and complete author derivation is retained.
\end{itemize}\medskip
\paragraph{2. Preliminaries: Notation.}
Throughout this paper, $[d]$ denotes the set $\{1,2,\ldots,d\}$. For a random vector (or bit-string) $z$ in $\mathbb R^n$, we will use $z_i$ or $z(i)$ to denote the $i$-th coordinate of $z$, depending on whether we need to use the subscript for another index. We shall use $\mathbf i=(i_1,\ldots,i_d)$ for a $d$-tuple of indices. For a $d$-tuple $\mathbf i$, we write $|\mathbf i|=d$ to denote the size of the tuple.
For a matrix $M\in\mathbb C^{N\times N}$, we denote by $M^*$ its conjugate transpose. Given a string $x\in\mathbb R^N$, the $N\times N$ diagonal matrix with $x$ on the diagonal is denoted by $\operatorname{Diag}(x)$. The normalized trace of an $N\times N$ matrix $M$ is defined as $\operatorname{tr}_N(M)=\frac1N(\sum_{i=1}^NM_{ii})$. The operator norm of a matrix $M$ is denoted by $\|M\|_{\mathrm{op}}$. The left (resp. right) polar decomposition $(V,P)$ of a square matrix $M$ is a factorization of the form $M=VP$ (resp. $M=PV$) where $V$ is a unitary matrix and $P$ is a positive semidefinite matrix --- such a factorization always exists for any square matrix, and can be obtained easily from the singular-value decomposition of $M$. An $N\times N$ matrix $U$ is called a Haar random unitary if it is distributed according to the Haar measure on the Unitary group $U(N)$.
Random variables are typically denoted by capital letters (e.g., $X$). We write $\mathbb E[f(X)]$ and $\operatorname{Var}[f(X)]$ to denote the expectation and variance of the random variable $f(X)$; if $f:\{\pm1\}^m\to\mathbb R$ then we abbreviate it to $\mathbb Ef$ and $\operatorname{Var}[f]$, where the expectation and variance are taken with respect to the uniform measure on the discrete cube $\{\pm1\}^m$.
\paragraph{Fourier Analysis on the Discrete Cube.}
We give some basic facts about Fourier analysis on the discrete cube and refer to the book [28] for more details. Every function $f:\{\pm1\}^m\to\mathbb R$ can be written uniquely as a sum of monomials $\chi_S(x)=\prod_{i\in S}x_i$,
\[f(x)=\sum_{S\subseteq[m]}\widehat f(S)\chi_S(x),\tag{5}\]
where $\widehat f(S)=\mathbb E[f(X)\chi_S(X)]$ is the Fourier coefficient with respect to the uniform $X\in\{\pm1\}^m$. The monomials $\chi_S(x)=\prod_{i\in S}x_i$ form an orthonormal basis for real-valued functions on $\{\pm1\}^m$, called the Fourier basis. Parseval's identity implies that for uniform $X\in\{\pm1\}^m$,
\[\mathbb Ef^2=\sum_{S\subseteq[m]}\widehat f(S)^2,\qquad\operatorname{Var}[f]=\sum_{S\subseteq[m]:S\neq\varnothing}\widehat f(S)^2.\]
For a function on the hypercube, we define $\|f\|_2^2:=\mathbb Ef^2$ which can also be viewed as the sum of squared Fourier coefficients because of Parseval's identity.
The discrete derivative of a function on the hypercube $\{\pm1\}^m$ is given by
\[\partial_i f(x)=\frac12(f(x^{i\to1})-f(x^{i\to-1})),\]
where $x^{i\to b}$ is the same as $x$ except that the $i$-th coordinate is set to $b$. It is easily checked that $\partial_i f(x)$ coincides with the real partial derivative $\partial/\partial x_i$ of the real multilinear polynomial given by (5).
For a real-valued function $f:\{\pm1\}^m\to\mathbb R$, the influence of a variable $x_i$ on $f$ is defined as
\[\operatorname{Inf}_i(f)=\mathbb E|\partial_i f|^2=\sum_{S\subseteq[m]:i\in S}\widehat f(S)^2.\]
\paragraph{Block-multilinear Forms.}
A degree-$d$ block-multilinear form $f:\{\pm1\}^{n\times d}\to\mathbb R$ is given by
\[f(\mathbf x_1,\ldots,\mathbf x_d)=\mathbb Ef+\sum_{m=1}^d\sum_{\substack{|\mathbf b|=m\\|\mathbf i|=m}}\widehat f_{\mathbf b,\mathbf i}\,x_{b_1}(i_1)x_{b_2}(i_2)\cdots x_{b_m}(i_m),\tag{6}\]
where the tuple $\mathbf b=(b_1,\ldots,b_m)$ satisfies $1\leq b_1<\cdots<b_m\leq d$ and $\mathbf i\in[n]^m$ is an $m$-tuple. The expectation that yields the constant term is uniform over $\{\pm1\}^{n\times d}$. Note that $m$ is determined from the size of the tuple $\mathbf b$, so we just write $\widehat f_{\mathbf b,\mathbf i}$ above.
From Parseval's identity, the variance of $f$ and the influence of $x_b(i)$ on $f$ (where $b\in[d]$ and $i\in[n]$) are respectively given by
\[\operatorname{Var}[f]=\sum_{m=1}^d\sum_{\substack{|\mathbf b|=m\\|\mathbf i|=m}}\widehat f_{\mathbf b,\mathbf i}^{\,2},\qquad\operatorname{Inf}_{b,i}(f)=\sum_{m=1}^d\sum_{\substack{|\mathbf b|=m:b\in\mathbf b\\|\mathbf i|=m:i_b=i}}\widehat f_{\mathbf b,\mathbf i}^{\,2}.\]
From the above, it follows that for any block $b\in[d]$,
\[\sum_{i\in[n]}\operatorname{Inf}_{b,i}(f)^2\leq\operatorname{Var}[f]\leq\sum_{b\in[d],i\in[n]}\operatorname{Inf}_{b,i}(f)^2,\tag{7}\]
where the first inequality is an equality if $f$ is a homogeneous degree-$d$ block-multilinear form. For any $b\in[d]$, we write $\operatorname{MaxInf}_b(T)=\max\{\operatorname{Inf}_{b,i}(T):i\in[n]\}$ to denote the maximum influence of any variable in the block $\mathbf x_b$.
Since each non-constant monomial contains at most one variable from each block and the ordering of the blocks is fixed, a degree-$d$ block-multilinear form $f:\{\pm1\}^{n\times d}\to\mathbb R$ can be naturally viewed as a non-commutative polynomial in matrix variables with the constant term replaced with $\mathbb Ef$ times the identity. Denoting the non-commutative polynomial as $f(\mathbf U_1,\ldots,\mathbf U_d)$ where each $\mathbf U_b=(U_b(1),\ldots,U_b(n))$ is a block of non-commutative variables, the completely bounded norm $\|f\|_{\mathrm{cb}}$ of the form $f$ is defined as
\begin{multline*}
\|f\|_{\mathrm{cb}}=\sup\{\|f(\mathbf U_1,\ldots,\mathbf U_d)\|_{\mathrm{op}}:N\in\mathbb N,\ U_b(i)\in\mathbb C^{N\times N},\\
\|U_b(i)\|_{\mathrm{op}}\leq1,\ b\in[d],\ i\in[n]\}.
\end{multline*}
The completely bounded norm originates in the theory of operator algebras. In the literature, this norm is sometimes defined for homogeneous block-multilinear forms only, but here we extend the definition to non-homogeneous block-multilinear forms.
The supremum above is always attained and can be computed by solving a semidefinite program as shown by Gribling and Laurent [20]. One can also equivalently restrict the supremum in the definition above to unitary matrices since the convex hull of the set of unitary matrices is the unit operator-norm ball. Moreover, $\|T\|_\infty\leq\|T\|_{\mathrm{cb}}$ where $\|T\|_\infty=\max_{x\in\{\pm1\}^{n\times d}}|T(x)|$, so forms that are completely bounded are also bounded on the hypercube.
Note that if $f$ is a degree-$d$ block-multilinear form and if we fix some of the input bits to $\pm1$, then the resulting function $g$ is also a degree-$d$ block-multilinear form with the same blocks $\mathbf x_1,\ldots,\mathbf x_d$, but it does not depend on the variables that were fixed. It is also easy to see that $\|g\|_{\mathrm{cb}}\leq\|f\|_{\mathrm{cb}}$ because while computing $\|g\|_{\mathrm{cb}}$ we may restrict the matrix variables to $\pm I$ if that particular variable was set to $\pm1$. In other words, completely bounded norm does not increase under restrictions.
We write $\operatorname{MaxInf}(f)=\max_{b\in[d],i\in[n]}\operatorname{Inf}_{b,i}(f)$ for the maximum influence of any variable in $f$.

\paragraph{Appendix A.1: Necessary free-probability definitions.}
There are many excellent books on free probability theory. In particular, we refer to the book [26] for more details than the brief introduction given here.
\paragraph{$C^*$-algebras.}
Let $\mathcal A$ be a unital $C^*$-algebra. For our purposes, we can think of this as an algebra of bounded operators on a complex Hilbert space which is self-adjoint ($a\in\mathcal A$ implies $a^*\in\mathcal A$), closed in the operator norm $\|\cdot\|$, and contains the identity ($1\in\mathcal A$). A faithful trace $\varphi$ on $\mathcal A$ is a continuous linear functional $\varphi:\mathcal A\to\mathbb C$ that is unital ($\varphi(1)=1$), positive $\varphi(aa^*)\geq0$, and $\varphi(aa^*)=0$ iff $a=0$.
The pair $(\mathcal A,\varphi)$ where $\mathcal A$ is a unital $C^*$-algebra and $\varphi$ is a faithful trace on $\mathcal A$ is called a $C^*$-probability space. Elements of $\mathcal A$ are called non-commutative random variables. An example of a finite-dimensional $C^*$-probability space is the class $(M_n(\mathbb C),\operatorname{tr}_n)$, which is the class of $n\times n$ complex matrices with the normalized trace functional defined as $\operatorname{tr}_n(M)=\frac1n\sum_{i=1}^nM_{ii}$. General $C^*$-probability spaces allow us to extend these definitions to infinite-dimensional operators, which are needed to define a non-commutative analog of independence called free independence. Faithfulness of the trace $\varphi$ then ensures that $\|a\|=\lim_{m\to\infty}\varphi((aa^*)^m)^{1/2m}$ (see [26, Proposition 3.17]). In particular, this allows one to compute the norm $\|\cdot\|$ by using the trace method and taking higher powers of the trace functional $\varphi$, as we will see later.
An illustrative example of an infinite-dimensional $C^*$-algebra is the following: let $G$ be an infinite discrete group with the identity element 1 and let $\ell_2(G)$ be the associated complex Hilbert space. The canonical orthonormal basis for $\ell_2(G)$ is given by the standard basis vectors $\{e_g:g\in G\}$. The left-regular representation of $G$ maps elements of $G$ to bounded linear operators on $\ell_2(G)$ and is defined by its action on the standard basis vectors,
\[\rho(g)e_h=e_{gh},\quad\forall h\in G.\]
The operator norm closure of the linear span of $\{\rho(g):g\in G\}$ forms a unital algebra $\mathcal A$ (called the reduced $C^*$-algebra of $G$) and together with the faithful trace $\varphi$ defined as
\[\varphi\left(\sum_g\alpha_g\rho(g)\right)=\alpha_1\]
gives a $C^*$-probability space.
\paragraph{Free Independence.}
Let $(\mathcal A,\varphi)$ be a $C^*$-probability space and let $\{\mathcal A_i\}_{i=1}^n$ be unital $*$-subalgebras of $\mathcal A$. They are said to be free (or freely independent) if for all $k\in[n]$, for all indices $i_1,\ldots,i_k\in[n]$, and for all $a_1\in\mathcal A_{i_1},\ldots,a_k\in\mathcal A_{i_k}$ satisfying $\varphi(a_1)=\cdots=\varphi(a_k)=0$, the joint free moment,
\[\varphi(a_1\cdots a_k)=0\]
whenever $j_1\neq j_2,j_2\neq j_3,\ldots,j_{k-1}\neq j_k$, that is, the free moments vanish when all the neighboring elements in the sequence $a_1,\ldots,a_k$ come from subalgebras with distinct indices, for example, $\varphi(a_1a_2a_1^*a_2^*a_3a_2)=0$.
Non-commutative random variables $a_1,\ldots,a_n\in(\mathcal A,\varphi)$ are said to be free if the subalgebras $\{\mathcal A_i\}_{i=1}^n$ are free, where $\mathcal A_i$ is the unital $*$-subalgebra generated by $a_i$ (the linear span of all monomials $a_i^{\epsilon_1}a_i^{\epsilon_2}\cdots a_i^{\epsilon_r}$ where $\epsilon_1,\ldots,\epsilon_r\in\{1,*\}$ and $r\in\mathbb N\cup\{0\}$). Note that the corresponding unital $C^*$-subalgebras obtained by taking the norm closure of each $\mathcal A_i$ are also freely independent in this case (see [26, Exercise 5.23]).
We remark that the set of free non-commutative random variables is an empty set if the underlying $C^*$-probability space is finite (for instance $(M_n(\mathbb C),\operatorname{tr}_n)$), so to find non-trivial examples one needs to work with infinite-dimensional $C^*$-probability spaces.
\paragraph{Free Haar Unitaries and Free Groups.}
Let $(\mathcal A,\varphi)$ be a $C^*$-probability space. An element $u\in\mathcal A$ is a Haar unitary if it is a unitary, i.e. $uu^*=u^*u=1$, and if $\varphi(u^k)=0$ for all non-zero integers $k$. A family $S=\{u_1,\ldots,u_n\}\in\mathcal A$ in a $C^*$-probability space $(\mathcal A,\varphi)$ is called a free Haar unitary family if each $u\in S$ is a Haar unitary and if $u_1,\ldots,u_n$ are free. For notational convenience, let us define $S^*=\{u_1^*,\ldots,u_n^*\}$ to be the set of corresponding adjoints.
One can give a concrete construction of free Haar unitaries in terms of the free group. The free group $F_n$ with a generating set $S$ of size $n$ is an infinite discrete group constructed as follows: a word is defined to be product of elements of $S\cup S^*$ with $\bot$ denoting the empty word that contains no symbols. A word is called reduced if it does not contain a sub-word of the form $gg^*$ or $g^*g$ for $g\in S$. Given a word that is not reduced, the process of repeatedly removing such sub-words until it becomes reduced is called reduction. The free group $F_n$ consists of all reduced words that can be built from the symbols in $S\cup S^*$ with the group operation being a product of words followed by reduction. The identity is the empty word $\bot$.
Let $\rho$ denote the left-regular representation of the free group $F_n$. Then, writing $S=\{s_1,\cdots,s_n\}$, the family $u_i=\rho(s_i)$ for $i\in[n]$, is a free Haar unitary family in the $C^*$-probability space given by the reduced $C^*$-algebra of the group $F_n$. This connection to the free group also allows us to give a very precise condition when the trace $\varphi$ evaluated on a non-commutative monomial in the $u_i$'s vanishes. For a $d$-tuple $\mathbf i=(i_1,\ldots,i_d)\in[m]^d$, let $u_{\mathbf i}$ denote the non-commutative monomial $u_{i_1}\cdots u_{i_d}$ and write $u_{\mathbf i}^*=(u_{\mathbf i})^*=u_{i_d}^*\cdots u_{i_1}^*$. Let $\mathbf i^1,\ldots,\mathbf i^t,\mathbf j^1,\ldots,\mathbf j^t$ each be a $d$-tuple in $[m]^d$ and consider the degree-$2td$ non-commutative monomial $w=u_{\mathbf i^1}u_{\mathbf j^1}^*u_{\mathbf i^2}u_{\mathbf j^2}^*\cdots u_{\mathbf i^t}u_{\mathbf j^t}^*$. Note that a degree-$2td$ monomial $w$ corresponds to an ordered $2td$-tuple of variables. To illustrate, if $t=1,m=3$ and $\mathbf i^1=(1,2,3)$ and $\mathbf j^1=(2,2,1)$, then $w=u_1u_2u_3(u_2u_2u_1)^*=u_1u_2u_3u_1^*u_2^*u_2^*$ and corresponds to the ordered tuple $(u_1,u_2,u_3,u_1^*,u_2^*,u_2^*)$. We can also interpret $w$ as a word in the free group by applying the reduction rules. Then the next proposition follows from the definitions of free independence and Haar unitaries.
\paragraph{Proposition 13.}
$\varphi(w)=1$ iff $w$ reduces to identity in the free group $F_n$, and $\varphi(w)=0$ otherwise.

\paragraph{Theorem 3.}
Let $f$ be a degree-$d$ block-multilinear form with $\|f\|_{\mathrm{cb}}\leq1$. Then, we have
\[\operatorname{MaxInf}(f)\geq\frac{(\operatorname{Var}[f])^2}{e(d+1)^4}.\]
\paragraph{Corollary 5.}
Let $\epsilon,\delta>0$ and let $f:\{\pm1\}^{n\times d}\to\mathbb R$ be a degree-$d$ block-multilinear form with $\|f\|_{\mathrm{cb}}\leq1$. Then, there is a deterministic classical algorithm that makes $O(d^5\epsilon^{-8}\delta^{-5})$ queries and approximates $f(x)$ up to an additive error $\epsilon$ on $1-\delta$ fraction of the inputs $x\in\{\pm1\}^{n\times d}$.

\paragraph{3.1. Low-degree polynomials of Haar random unitaries.}
As discussed in the proof overview, we require bounds on the operator norm (as well as normalized trace) of low-degree polynomials of random unitaries and these follow from known results in free probability theory. Here we explain these connections and also prove some auxillary lemmas needed for the proof of Theorem 4 and Theorem 3.
Let $z_{\mathbf i}$ denote the non-commutative monomial $z_{i_1}z_{i_2}\cdots z_{i_d}$ for a $d$-tuple $\mathbf i=(i_1,\ldots,i_d)\in[t]^d$ and let $p(z_1,\ldots,z_t)$ be a non-commutative polynomial in the variables $z_1,\ldots,z_t$. We are interested in computing the operator norm $\|\cdot\|_{\mathrm{op}}$ and the normalized trace $\operatorname{tr}_N$ of the polynomial $p(z_1,\ldots,z_t)$ (or its higher moments) when substituting $N\times N$ Haar random unitaries for the variables $z_i$.
As explained previously, the theory of free probability gives us tools that allow us to compute the above in the limit $N\to\infty$. In particular, Voiculescu [35] showed that the (normalized) trace of polynomials in Haar random unitaries and their conjugates converges to the trace of the same polynomial evaluated on certain infinite-dimensional operators called Haar unitaries that satisfy a non-commutative notion of independence called free independence. This was strengthened by Collins and Male [16] who showed that such convergence also holds for the operator norm. A short primer on free probability is given in Appendix A.1, but for now one can think of $\mathcal A$ as a self-adjoint algebra of bounded linear operators on a Hilbert space and $\varphi$ as a trace functional for such operators in the statement given below.
\paragraph{Theorem 8 ([35,16]).}
Let $p(z_1,\ldots,z_{2t})$ be a non-commutative polynomial in $\mathbb R\langle z_1,\ldots,z_{2t}\rangle$. If $U_1,\ldots,U_t$ are $N\times N$ Haar random unitaries, then almost surely,
\[\lim_{N\to\infty}\operatorname{tr}_N[p(U_1,\ldots,U_t,U_1^*,\ldots,U_t^*)]=\varphi[p(u_1,\ldots,u_t,u_1^*,\ldots,u_t^*)],\]
\[\lim_{N\to\infty}\|p(U_1,\ldots,U_t,U_1^*,\ldots,U_t^*)\|_{\mathrm{op}}=\|p(u_1,\ldots,u_t,u_1^*,\ldots,u_t^*)\|,\]
where $u_1,\ldots,u_t$ are free Haar unitaries in a $C^*$-probability space $(\mathcal A,\varphi)$ and $\|\cdot\|$ is the norm for the underlying $C^*$-algebra.
Using the above result it suffices to consider free Haar unitaries in a $C^*$-probability space to compute the operator norm and trace of polynomials of random unitaries. For a non-commutative polynomial $p(z_1,\ldots,z_t)=\sum_{|\mathbf i|\leq d}c_{\mathbf i}z_{\mathbf i}$, denoting by $\|p\|_2=(\sum_{|\mathbf i|\leq d}|c_{\mathbf i}|^2)^{1/2}$, one can show the following easily using techniques from free probability.
\paragraph{Lemma 9.}
Let $p(z_1,\ldots,z_t)=\sum_{|\mathbf i|\leq d}c_{\mathbf i}z_{\mathbf i}$ be a non-commutative degree-$d$ polynomial in $\mathbb R\langle z_1,\ldots,z_t\rangle$ and $u_1,\ldots,u_t$ be free Haar unitaries in a $C^*$-probability space $(\mathcal A,\varphi)$. Then,
\[\varphi[p(u_1,\ldots,u_t)(p(u_1,\ldots,u_t))^*]=\|p\|_2^2.\]
\begin{proof}[Proof of Lemma 9 (Appendix A.2)]
Writing $u_{\mathbf i}^*=(u_{\mathbf i})^*$ for a tuple $\mathbf i$ and using linearity of $\varphi$, we have that
\[\varphi[p(u_1,\ldots,u_t)(p(u_1,\ldots,u_t))^*]=\sum_{|\mathbf i|,|\mathbf j|\leq d}c_{\mathbf i}c_{\mathbf j}\varphi(u_{\mathbf i}u_{\mathbf j}^*).\]
From Proposition 13, the term $\varphi(u_{\mathbf i}u_{\mathbf j}^*)$ is 1 iff $u_{\mathbf i}u_{\mathbf j}^*$ reduces to identity in the free group $F_t$ with generators $u_1,\ldots,u_t$. For the right-hand side above, this only happens when $\mathbf i=\mathbf j$ and thus these are the only non-zero terms. Thus,
\[\varphi[p(u_1,\ldots,u_t)(p(u_1,\ldots,u_t))^*]=\sum_{|\mathbf i|\leq d}|c_{\mathbf i}|^2.\]
\end{proof}
The above implies that $\operatorname{tr}_N[p(U_1,\ldots,U_t)(p(U_1,\ldots,U_t))^*]$ converges to $\|p\|_2^2$ almost surely as $N\to\infty$.
In particular, we will rely on the following result that follows from the work of Kemp and Speicher [23] who consider the operator norm of homogeneous polynomials evaluated on free $R$-diagonal operators, a class that includes free Haar unitaries as well. We also remark that a bound where the right-hand side below is worse by a multiplicative $O(d^{1/2})$ factor also follows from the work of Haagerup [22] who proved it in another context, predating even the introduction of free probability theory. Haagerup considered the more general case of polynomials in both $u_i$'s and $u_i^*$'s.
\paragraph{Theorem 10 ([23]).}
Let $p(z_1,\ldots,z_t)=\sum_{|\mathbf i|=d}c_{\mathbf i}z_{\mathbf i}$ be a homogeneous non-commutative degree-$d$ polynomial in $\mathbb R\langle z_1,\ldots,z_t\rangle$ and $u_1,\ldots,u_t$ be free Haar unitaries in a $C^*$-probability space. Then,
\[\|p(u_1,\ldots,u_t)\|\leq\sqrt{e(d+1)}\cdot\|p\|_2,\]
where the left-hand side denotes the norm in the underlying $C^*$-algebra.
We shall need to extend the above bound to non-homogeneous polynomials. Let $p(z_1,\ldots,z_t)=\sum_{|\mathbf i|\leq d}c_{\mathbf i}z_{\mathbf i}$ and let $p_k(z_1,\ldots,z_t)=\sum_{|\mathbf i|=k}c_{\mathbf i}z_{\mathbf i}$ denote the degree-$k$ homogeneous part of $p$. Writing $p_k=p_k(u_1,\ldots,u_t)$ for $0\leq k\leq d$ and $p=p(u_1,\ldots,u_t)$, it follows from the triangle inequality, Theorem 10, and Cauchy--Schwarz, that
\begin{align*}
\|p\|&\leq\sum_{k=0}^d\|p_k\|\leq\sum_{k=0}^d\sqrt{e(k+1)}\|p_k\|_2\\
&\leq\sqrt e\left(\sum_{k=0}^d(k+1)\right)^{1/2}\left(\sum_{k=0}^d\|p_k\|_2^2\right)^{1/2}\leq\sqrt e(d+1)\cdot\|p\|_2.
\end{align*}
Thus, we essentially get the same bound as in the homogeneous case, at the expense of an additional $O(d^{1/2})$ factor.
Collecting all the above we have the following as a direct consequence:
\paragraph{Theorem 11.}
Let $p(z_1,\ldots,z_t)=\sum_{|\mathbf i|\leq d}c_{\mathbf i}z_{\mathbf i}$ be a non-commutative degree-$d$ polynomial in $\mathbb R\langle z_1,\ldots,z_t\rangle$ and $U_1,\ldots,U_t$ be independent $N\times N$ Haar random unitaries. Then, the following holds almost surely,
\[\lim_{N\to\infty}\operatorname{tr}_N[p(U_1,\ldots,U_t)(p(U_1,\ldots,U_t))^*]=\|p\|_2^2,\]
and
\[\lim_{N\to\infty}\|p(U_1,\ldots,U_t)\|_{\mathrm{op}}\leq\sqrt e(d+1)\cdot\|p\|_2.\]
Moreover, the factor $(d+1)$ in the operator norm bound can be improved to $\sqrt{d+1}$ if additionally $p$ is also assumed to be homogeneous.
Based on the above theorem, we prove the following key lemma which captures the polar decomposition strategy mentioned in the earlier proof overview (Section 1.2). This will serve as the key ingredient in the proof of Theorem 3 and Theorem 4.
\paragraph{Lemma 12.}
Let $p$ be a non-commutative degree-$d$ polynomial in $\mathbb R\langle y_1,\ldots,y_m,z_1,\ldots,z_t\rangle$ given by
\[p(y_1,\ldots,y_m,z_1,\ldots,z_t)=\sum_{i=1}^my_iq_i(z_1,\ldots,z_t)+q_0(z_1,\ldots,z_t).\]
Then, for every $\delta>0$, there exist an integer $N$ and $N\times N$ unitaries $V_1,\ldots,V_m,W_1,\ldots,W_t$ such that
\[\|p(V_1,\ldots,V_m,W_1,\ldots,W_t)\|_{\mathrm{op}}\geq\frac1{\sqrt e(d+1)}\sum_{i=1}^m\|q_i\|_2-\delta.\]
Moreover, the factor in front can be improved to $(e(d+1))^{-1/2}$ if $p$ is homogeneous as well.
We note that the definition of completely bounded norm allows us to plug in matrices of arbitrarily large dimensions, so the exact dependence on $\delta$ is not relevant for our purposes since we can always make it arbitrarily small.
\begin{proof}
For an arbitrary integer $N$, let us pick independent $N\times N$ Haar random unitaries $W_1,\ldots,W_t$ which we substitute for the variables $z_1,\ldots,z_t$, respectively, and let $M_i=q_i(W_1,\ldots,W_t)$ be the corresponding random matrices. Then, for any tuple of matrices $V_1,\ldots,V_m$ that we could substitute for the variables $y_1,\ldots,y_m$, we have that
\[p(V_1,\ldots,V_m,W_1,\ldots,W_t)=\sum_{i=1}^mV_iM_i+M_0.\]
Theorem 11 and union bound imply that as $N\to\infty$, with probability 1 all the following events simultaneously hold:
\begin{itemize}
\item $\|M_i\|_{\mathrm{op}}\leq\sqrt e(d+1)\cdot\|q_i\|_2$ for each $i$,
\item $\operatorname{tr}_N(M_i^*M_i)=\|q_i\|_2^2$ for each $i$, where $\operatorname{tr}_N(M)$ is the normalized trace.
\end{itemize}
To show that the operator norm must be large, let us fix a sufficiently large $N$ and a choice of $N\times N$ unitaries $W_1,\ldots,W_t$ such that $M_i$ satisfies $\|M_i\|_{\mathrm{op}}\leq\sqrt e(d+1)\cdot\|q_i\|_2+\epsilon$ and $\operatorname{tr}_N(M_i^*M_i)\geq\|q_i\|_2^2-\epsilon$ for each $0\leq i\leq m$, where $\epsilon$ can be made arbitrarily small by increasing $N$. For $0\leq i\leq m$, let $M_i=U_iP_i$ be the left polar decomposition of $M_i$, where $U_i$ is a unitary matrix and $P_i$ is a positive semidefinite matrix.
We select the tuple of unitary matrices $V_1,\ldots,V_m$ that we substitute for the variables $y_1,\ldots,y_m$ to be $V_i=U_0U_i^*$ for $i\in[m]$. With this we have that $\|p(V_1,\ldots,V_m,W_1,\ldots,W_t)\|_{\mathrm{op}}$ is at least
\begin{align*}
\left\|M_0+\sum_{i=1}^mV_iM_i\right\|_{\mathrm{op}}
&=\left\|U_0P_0+\sum_{i=1}^mU_0U_i^*U_iP_i\right\|_{\mathrm{op}}\\
&=\left\|U_0P_0+\sum_{i=1}^mU_0P_i\right\|_{\mathrm{op}}\\
&=\left\|P_0+\sum_{i=1}^mP_i\right\|_{\mathrm{op}}\\
&\geq\operatorname{tr}_N\left(P_0+\sum_{i=1}^mP_i\right)\geq\operatorname{tr}_N\left(\sum_{i=1}^mP_i\right),
\end{align*}
where the last equality follows since the operator norm is unitarily invariant and the last two inequalities follow from the positive semidefiniteness of the $P_i$'s.
For every positive semidefinite matrix $P$, we have that $\operatorname{tr}_N(P)\geq\operatorname{tr}_N(P^2)/\|P\|_{\mathrm{op}}$. Hence,
\[\|p(V_1,\ldots,V_m,W_1,\ldots,W_t)\|_{\mathrm{op}}\geq\sum_{i=1}^m\frac{\operatorname{tr}_N(P_i^2)}{\|P_i\|_{\mathrm{op}}}.\]
By our choice of $M_i$, we have that $\operatorname{tr}_N(P_i^2)=\operatorname{tr}_N(M_i^*M_i)\geq\|q_i\|_2^2-\epsilon$ and $\|P_i\|_{\mathrm{op}}=\|M_i\|_{\mathrm{op}}\leq\sqrt e(d+1)\|q_i\|_2+\epsilon$. Since $\epsilon$ can be made arbitrarily small by increasing $N$, it follows that
\[\|p(V_1,\ldots,V_m,W_1,\ldots,W_t)\|_{\mathrm{op}}\geq\frac1{\sqrt e(d+1)}\sum_{i=1}^m\|q_i\|_2-\delta,\]
for large enough $N$. The improved bound for the homogeneous case follows directly by plugging the bound of Theorem 11 into the above proof.
\end{proof}

\paragraph{3.2. Non-commutative root-influence inequality.}
For clarity in the proofs below, we remind our convention that all tuples or blocks are denoted with boldface fonts (e.g. $\mathbf U_1$ or $\mathbf A$), while a single element is denoted without boldface (e.g. $U_1(i)$ or $A_i$ or $A$). Before proceeding with the proof, we restate the statement for convenience.
\paragraph{Theorem 4 (Non-commutative root-influence inequality).}
Let $f$ be a homogeneous degree-$d$ block-multilinear form. Then, for blocks $b\in\{1,d\}$,
\[\|f\|_{\mathrm{cb}}\geq\frac1{\sqrt{e(d+1)}}\sum_{i=1}^n\sqrt{\operatorname{Inf}_{b,i}(f)}.\]
\begin{proof}[Proof of Theorem 4]
Since $f$ is homogeneous, we can write
\begin{align*}
f(\mathbf x_1,\ldots,\mathbf x_d)
&=\sum_{i_1,\ldots,i_d\in[n]}\widehat f_{i_1,\ldots,i_d}\,x_1(i_1)x_2(i_2)\cdots x_d(i_d)\\
&=\sum_{i=1}^nx_1(i)\underbrace{\left(\sum_{i_2,\ldots,i_d\in[n]}\widehat f_{i,\ldots,i_d}\,x_2(i_2)\cdots x_d(i_d)\right)}_{:=f_i(\mathbf x_2,\ldots,\mathbf x_d)}.
\end{align*}
In this case, it follows from (7) that for each $i\in[n]$, we have
\[\operatorname{Var}[f_i]=\|f_i\|_2^2=\operatorname{Inf}_{1,i}(f)\quad\text{and}\quad\operatorname{Var}[f]=\sum_{i=1}^n\operatorname{Inf}_{1,i}(f).\tag{8}\]
Let us denote the corresponding non-commutative block-multilinear polynomials by $f(\mathbf U_1,\ldots,\mathbf U_d)$ and $f_i(\mathbf U_2,\ldots,\mathbf U_d)$ where $\mathbf U_b=(U_b(1),\ldots,U_b(n))$ denotes the $b$th block of non-commutative variables. To show a lower bound on $\|f\|_{\mathrm{cb}}$ it suffices to exhibit a collection of square matrices $\{U_b(i)\}_{b\in[d],i\in[n]}$ with operator norm at most 1, such that $\|f(\mathbf U_1,\ldots,\mathbf U_d)\|_{\mathrm{op}}$ is large.
Applying Lemma 12 for the homogeneous case (with $p=f$, $q_i=f_i$ for $i\in[n]$, and $q_0=0$), it follows that for every $\delta>0$ there exists an integer $N$ and a choice of tuples $\mathbf U_1,\ldots,\mathbf U_d$ of $N\times N$ unitaries such that
\begin{align*}
\|f\|_{\mathrm{cb}}&\geq\|f(\mathbf U_1,\ldots,\mathbf U_d)\|_{\mathrm{op}}\\
&\geq\frac1{\sqrt{e(d+1)}}\sum_{i\in[n]}\|f_i\|_2-\delta\\
&\overset{(8)}=\frac1{\sqrt{e(d+1)}}\left(\sum_{i=1}^n\sqrt{\operatorname{Inf}_{1,i}(f)}\right)-\delta.
\end{align*}
Taking $\delta\to0$, we get the statement of the lemma. The proof for the inequality when $b=d$ is the last block follows similarly by using the right polar decomposition analogue of Lemma 12.
\end{proof}
\paragraph{3.3. Aaronson--Ambainis Conjecture for non-homogeneous forms.}
In this section, we prove Theorem 3, which requires handling non-homogeneous forms. The proof will be similar to the proof of Theorem 4 but we will need to be careful about certain details.
\begin{proof}[Proof of Theorem 3]
Any block-multilinear polynomial $f(\mathbf x_1,\ldots,\mathbf x_d)$ can be written as
\[f(\mathbf x_1,\ldots,\mathbf x_d)=\mathbb Ef+\sum_{b\in[d]}f_b(\mathbf x_b,\mathbf x_{b+1},\ldots,\mathbf x_d),\]
where $f_b$ consists of all monomials of $f$ that start with a variable in the $b$th block $\mathbf x_b$. Note that $f_b$ depends only on the variables in blocks $\mathbf x_b,\mathbf x_{b+1},\ldots,\mathbf x_d$. Moreover, it follows from Parseval that
\[\operatorname{Var}[f]=\sum_{b\in[d]}\|f_b\|_2^2=\sum_{b\in[d]}\operatorname{Var}[f_b],\tag{9}\]
so there exists a block $\beta\in[d]$ such that $\operatorname{Var}[f_\beta]\geq\frac1d\operatorname{Var}[f]$.
Since $f_\beta$ contributes a lot to the variance, it is natural to try to find an influential variable in the block $\mathbf x_\beta$. Towards this end, we pull out the variables $x_\beta(i)$ and write
\[f_\beta(\mathbf x_\beta,\ldots,\mathbf x_d)=\sum_{i\in[n]}x_\beta(i)f_{\beta,i}(\mathbf x_{\beta+1},\ldots,\mathbf x_d),\]
for block-multilinear polynomials $f_{\beta,i}(\mathbf x_{\beta+1},\ldots,\mathbf x_d)$. Note that some of the $f_{\beta,i}$'s could be identically zero, so let us define $S$ to be the set of those $i$ such that $f_{\beta,i}$ is non-zero. We note that
\[\|f_{\beta,i}\|_2^2=\operatorname{Inf}_{\beta,i}(f_\beta)\leq\operatorname{Inf}_{\beta,i}(f),\tag{10}\]
which implies (using Parseval's identity) that
\[\frac1d\operatorname{Var}[f]\leq\operatorname{Var}[f_\beta]=\sum_{i\in S}\|f_{\beta,i}\|_2^2=\sum_{i\in S}\operatorname{Inf}_{\beta,i}(f_\beta).\tag{11}\]
Denote the corresponding non-commutative block-multilinear polynomials by $f(\mathbf U_1,\ldots,\mathbf U_d)$, $f_b(\mathbf U_b,\ldots,\mathbf U_d)$, and $f_\beta(\mathbf U_{\beta+1},\ldots,\mathbf U_d)$ where $\mathbf U_b=(U_b(1),\ldots,U_b(n))$ denotes the $b$th block of non-commutative variables. To show a lower bound on $\|f\|_{\mathrm{cb}}$ it suffices to exhibit a collection of square matrices $\{U_b(i)\}_{b\in[d],i\in[n]}$ with operator norm at most 1 such that $\|f(\mathbf U_1,\ldots,\mathbf U_d)\|_{\mathrm{op}}$ is large.
We set the matrices in blocks $\mathbf U_1,\ldots,\mathbf U_{\beta-1}$ to be zero (that is, the all-zero matrix $0$). Note that with this choice all polynomials $f_b(\mathbf U_b,\ldots,\mathbf U_d)$ where $b<\beta$ vanish and the non-commutative polynomial becomes
\begin{align*}
f(0,\ldots,0,\mathbf U_\beta,\mathbf U_{\beta+1},\ldots,\mathbf U_d)
&=\sum_{i\in S}U_\beta(i)f_{\beta,i}(\mathbf U_{\beta+1},\ldots,\mathbf U_d)\\
&\quad+\sum_{b=\beta+1}^df_b(\mathbf U_b,\mathbf U_{b+1},\ldots,\mathbf U_d)+\mathbb Ef,
\end{align*}
which is a non-commutative polynomial of the form considered in Lemma 12 (with $m=|S|$, $q_i=f_{\beta,i}$ and $q_0=\sum_{b=\beta+1}^df_b+\mathbb Ef$). Thus, by Lemma 12 for every small $\delta>0$ there exists an integer $N$ and a choice of $N\times N$ matrices for the blocks $\mathbf U_\beta,\ldots,\mathbf U_d$ such that
\begin{align*}
\|f\|_{\mathrm{cb}}&\geq\|f(0,\ldots,0,\mathbf U_\beta,\mathbf U_{\beta+1},\ldots,\mathbf U_d)\|_{\mathrm{op}}\\
&\geq\frac1{\sqrt e(d+1)}\sum_{i\in S}\|f_{\beta,i}\|_2-\delta\\
&\overset{(10)}=\frac1{\sqrt e(d+1)}\left(\sum_{i\in S}\sqrt{\operatorname{Inf}_{\beta,i}(f_\beta)}\right)-\delta\\
&\overset{(11)}\geq\frac1{\sqrt e(d+1)}\left(\frac{\sum_{i\in S}\operatorname{Inf}_{\beta,i}(f_\beta)}{\sqrt{\operatorname{MaxInf}(f)}}\right)-\delta\\
&\overset{(10)}\geq\frac1{\sqrt e(d+1)^2}\left(\frac{\operatorname{Var}[f]}{\sqrt{\operatorname{MaxInf}(f)}}\right)-\delta.
\end{align*}
Taking $\delta\to0$ and using the assumption that $\|f\|_{\mathrm{cb}}\leq1$, we obtain the statement of the theorem:
\[1\geq\|f\|_{\mathrm{cb}}\geq\frac1{\sqrt e(d+1)^2}\cdot\frac{\operatorname{Var}[f]}{\sqrt{\operatorname{MaxInf}(f)}}\quad\Longrightarrow\quad\operatorname{MaxInf}(f)\geq\frac{(\operatorname{Var}[f])^2}{e(d+1)^4}.\]
\end{proof}
\paragraph{3.4. Approximating completely bounded forms with decision trees.}
In this section, we briefly mention how to obtain Corollary 5. Aaronson and Ambainis [1, Theorem 3.3] showed that querying the most influential variable reduces the variance of the function $f$, and if that influence is lower bounded by a polynomial in $\operatorname{Var}[f]/d$, then after $\operatorname{poly}(d)$ queries (the exact quantitative dependence can be read off from their proof), the variance of the function becomes small enough so that it can be approximated almost-everywhere by its expectation. Since the family of degree-$d$ block-multilinear forms with completely bounded norm at most one is closed under restrictions, one can apply Theorem 3 repeatedly. This gives us Corollary 5.

\begin{thebibliography}{99}
\bibitem[1]{published-1} Scott Aaronson and Andris Ambainis. The need for structure in quantum speedups. Theory of Computing, 10(6):133–166, 2014.
\bibitem[2]{published-2} Scott Aaronson and Andris Ambainis. Forrelation: A problem that optimally separates quantum from classical computing. SIAM Journal on Computing, 47(3):982–1038, 2018.
\bibitem[3]{published-3} Scott Aaronson, Shalev Ben-David, and Robin Kothari. Separations in query complexity using cheat sheets. In Proceedings of the 48th Annual ACM Symposium on Theory of Computing, pages 863–876, 2016.
\bibitem[4]{published-4} Scott Aaronson, Shalev Ben-David, Robin Kothari, Shravas Rao, and Avishay Tal. Degree vs. approximate degree and quantum implications of Huang’s sensitivity theorem. In Proceedings of the 53rd Annual ACM Symposium on Theory of Computing, pages 1330–1342, 2021. arXiv: 2010.12629.
\bibitem[5]{published-5} Scott Aaronson, DeVon Ingram, and William Kretschmer. The acrobatics of BQP. CoRR, abs/2111.10409, 2021. arXiv:2111.10409.
\bibitem[6]{published-6} Leonard M. Adleman, Jonathan Demarrais, and Ming-Deh A. Huang. Quantum computability. SIAM Journal on Computing, 26(5):1524–1540, 1997.
\bibitem[7]{published-7} Andris Ambainis. Polynomial degree vs. quantum query complexity. Journal of Computer and System Sciences, 72(2):220–238, 2006.
\bibitem[8]{published-8} Andris Ambainis, Kaspars Balodis, Aleksandrs Belovs, Troy Lee, Miklos Santha, and Juris Smotrovs. Separations in query complexity based on pointer functions. Journal of the ACM, 64(5):32:1–32:24, 2017.
\bibitem[9]{published-9} Srinivasan Arunachalam, Jop Briët, and Carlos Palazuelos. Quantum query algorithms are completely bounded forms. SIAM Journal on Computing, 48(3):903–925, 2019.
\bibitem[10]{published-10} Nikhil Bansal and Makrand Sinha. k-Forrelation optimally separates quantum and classical query complexity. In Proceedings of the 53rd Annual ACM SIGACT Symposium on Theory of Computing, pages 1303–1316, 2021.
\bibitem[11]{published-11} Howard Barnum, Michael E. Saks, and Mario Szegedy. Quantum query complexity and semi-definite programming. In Proceedings of 18th Annual IEEE Conference on Computational Complexity, pages 179–193, 2003.
\bibitem[12]{published-12} Robert Beals, Harry Buhrman, Richard Cleve, Michele Mosca, and Ronald de Wolf. Quantum lower bounds by polynomials. Journal of the ACM, 48(4):778–797, 2001.
\bibitem[13]{published-13} Shalev Ben-David, Pooya Hatami, and Avishay Tal. Low-sensitivity functions from unam- biguous certificates. In Proceedings of the 8th Innovations in Theoretical Computer Science Conference, volume 67 of LIPIcs, pages 28:1–28:23. Schloss Dagstuhl - Leibniz-Zentrum für Informatik, 2017.
\bibitem[14]{published-14} Sergey Bravyi, David Gosset, Daniel Grier, and Luke Schaeffer. Classical algorithms for forrelation, 2021. arXiv:2102.06963.
\bibitem[15]{published-15} Harry Buhrman and Ronald de Wolf. Complexity measures and decision tree complexity: a survey. Theoretical Computer Science, 288(1):21–43, 2002.
\bibitem[16]{published-16} Benoît Collins and Camille Male. The strong asymptotic freeness of Haar and deterministic matrices. Annales Scientifiques de l’ENS, (4) 47, fascicule 1:147–163, 2014. arXiv:1105.4345.
\bibitem[17]{published-17} Benoît Collins and Ion Nechita. Random matrix techniques in quantum information theory. Journal of Mathematical Physics, 57(1):015215, 2016.
\bibitem[18]{published-18} Andreas Defant, Mieczyslaw Mastylo, and Antonio Pérez. On the Fourier spectrum of functions on Boolean cubes. Mathematische Annalen, 374:653–680, 2018.
\bibitem[19]{published-19} Irit Dinur, Ehud Friedgut, Guy Kindler, and Ryan O’Donnell. On the Fourier tails of bounded functions over the discrete cube. In Proceedings of the 38th Annual ACM Symposium on Theory of Computing, pages 437–446, 2006.
\bibitem[20]{published-20} Sander Gribling and Monique Laurent. Semidefinite programming formulations for the completely bounded norm of a tensor, 2019. arXiv:1901.04921.
\bibitem[21]{published-21} Lov K. Grover. A fast quantum mechanical algorithm for database search. In Proceedings of the 28th Annual ACM Symposium on Theory of Computing, pages 212–219, 1996.
\bibitem[22]{published-22} Uffe Haagerup. An example of a non nuclear C $\ast$ -algebra, which has the metric approximation property. Inventiones Mathematicae, 50:279–293, 1978/79.
\bibitem[23]{published-23} Todd Kemp and Roland Speicher. Strong Haagerup inequalities for free R-diagonal elements. Journal of Functional Analysis, 251(1):141–173, 2007. arXiv:math/0512481.
\bibitem[24]{published-24} Gatis Midrijanis. On randomized and quantum query complexities, 2005. arXiv:quant-ph/ 0501142.
\bibitem[25]{published-25} Ashley Montanaro. Some applications of hypercontractive inequalities in quantum information theory. Journal of Mathematical Physics, 53(12):122206, 2012.
\bibitem[26]{published-26} Alexandru Nica and Roland Speicher. Lectures on the Combinatorics of Free Probability. London Mathematical Society Lecture Note Series. Cambridge University Press, 2006.
\bibitem[27]{published-27} Noam Nisan and Mario Szegedy. On the degree of Boolean functions as real polynomials. Computational Complexity, 4(4):301–313, 1994.
\bibitem[28]{published-28} Ryan O’Donnell. Analysis of Boolean Functions. Cambridge University Press, 2014.
\bibitem[29]{published-29} Ryan O’Donnell, Michael E. Saks, Oded Schramm, and Rocco A. Servedio. Every decision tree has an influential variable. In Proceedings of the 46th Annual IEEE Symposium on Foundations of Computer Science, pages 31–39, 2005.
\bibitem[30]{published-30} Ryan O’Donnell and Yu Zhao. Polynomial bounds for decoupling, with applications. In Proceedings of 31st Conference on Computational Complexity, volume 50 of LIPIcs, pages 24:1–24:18. Schloss Dagstuhl - Leibniz-Zentrum für Informatik, 2016.
\bibitem[31]{published-31} Alexander A. Sherstov, Andrey A. Storozhenko, and Pei Wu. An optimal separation of randomized and quantum query complexity. In Proceedings of the 53rd Annual ACM SIGACT Symposium on Theory of Computing, pages 1289–1302, 2021.
\bibitem[32]{published-32} Peter W. Shor. Polynomial-time algorithms for prime factorization and discrete logarithms on a quantum computer. SIAM Journal on Computing, 26(5):1484–1509, 1997.
\bibitem[33]{published-33} Daniel R. Simon. On the power of quantum computation. SIAM Journal on Computing, 26(5):1474–1483, 1997.
\bibitem[34]{published-34} Avishay Tal. Towards optimal separations between quantum and randomized query complexities. In Proceedings of the 61st IEEE Annual Symposium on Foundations of Computer Science, pages 228–239, 2020.
\bibitem[35]{published-35} Dan Voiculescu. A strengthened asymptotic freeness result for random matrices with applica- tions to free entropy. International Mathematics Research Notices, 1998:41–63, 1998.
\bibitem[36]{published-36} Z. Yin, A. W. Harrow, M. Horodecki, M. Marciniak, and A. Rutkowski. Random and free observables saturate the Tsirelson bound for CHSH inequality. Physical Review A, 95(032101), 2017. arXiv:1512.00223.
\end{thebibliography}
\end{document}
